% =====================================================================
%  Chapitre 1 -- Du MLP aux CNN
%  Partie d : recapitulatif, quiz de synthese, TP1
% =====================================================================

\begin{frame}{MLP contre CNN : le tableau de bord}
\begin{center}
\small
\begin{tabular}{@{}p{0.30\textwidth}p{0.32\textwidth}p{0.32\textwidth}@{}}
\toprule
 & \textbf{MLP (couche dense)} & \textbf{CNN (couche convolutive)} \\
\midrule
\textbf{Entrée attendue} & vecteur (image aplatie) & tenseur (image, structure conservée) \\
\textbf{Voisinage} & tout l'image & un voisinage local ($3 \times 3$, $5 \times 5$) \\
\textbf{Poids} & un par (pixel, neurone) & un par position du noyau, partagés \\
\textbf{Paramètres, photo 24 Mpx} & $\approx 48$ milliards & $\approx 1\,800$ \\
\textbf{Dépend à la position} & oui & non (invariance par translation) \\
\textbf{Ce qu'il apprend} & des corrélations globales & des motifs locaux réutilisables \\
\textbf{Adapté aux images} & non, sauf très petites & oui \\
\bottomrule
\end{tabular}
\end{center}
\vspace{1mm}
\begin{kbox}[La phrase à retenir]
Un CNN n'est pas un MLP amélioré : c'est un MLP auquel on a \textbf{interdit}
de connecter des pixels éloignés, et \textbf{obligé} de réutiliser ses poids.
\end{kbox}
\end{frame}

\begin{frame}{À retenir}
\begin{kbox}[Les trois idées de ce chapitre]
\begin{enumerate}
  \item Une image est une \textbf{matrice} (niveaux de gris) ou un
  \textbf{tenseur} (couleur). L'aplatir détruit l'information la plus utile.
  \item Le MLP explose : le nombre de paramètres croît comme
  $n_{\text{pixels}} \times n_{\text{neurones}}$. Sur une vraie photo, c'est
  des dizaines de milliards.
  \item Deux idées règlent le problème : \textbf{localité} (un noyau ne voit
  qu'un voisinage) et \textbf{partage de poids} (le même noyau partout),
  d'où l'\textbf{invariance par translation}.
\end{enumerate}
\end{kbox}
\vspace{2mm}
\begin{wbox}[Ce qui reste à construire]
Nous avons le \emph{principe}. Il manque : le calcul précis de la convolution
(chapitre 2), l'architecture complète (chapitre 3), et l'entraînement
(chapitre 4).
\end{wbox}
\end{frame}

\begin{frame}{Quiz de synthèse du chapitre 1}
\begin{qbox}[Question 1]
Pourquoi aplatisse-t-on une image avant un MLP ?
\begin{itemize}
  \item[(a)] Pour réduire la mémoire \quad (b) Parce que l'entrée d'un MLP est un
  vecteur \quad (c) Pour normaliser les valeurs
\end{itemize}
\end{qbox}
\begin{qbox}[Question 2]
Une couche convolutive $5 \times 5$, $32$ filtres, sur une entrée à $3$ canaux
contient :
\begin{itemize}
  \item[(a)] $32$ paramètres \quad (b) $2\,432$ \quad (c) $4\,800$
\end{itemize}
\end{qbox}
\begin{qbox}[Question 3]
Lequel de ces défauts le partage de poids ne règle-t-il \textbf{pas} à lui seul ?
\begin{itemize}
  \item[(a)] L'explosion des paramètres \quad (b) La dépendance à la position
  \quad (c) Le fait qu'un filtre ne détecte qu'un motif
\end{itemize}
\end{qbox}
\begin{solbox}[Réponses]
\textbf{1.} (b). \quad
\textbf{2.} (b) : $5 \times 5 \times 3 \times 32 + 32 = 2\,400 + 32 = 2\,432$. \quad
\textbf{3.} (c) : il faut \emph{plusieurs} filtres, d'où les canaux de sortie
(chapitre 3).
\end{solbox}
\end{frame}

\begin{frame}{Travail pratique 1 --- noyaux de convolution}
\begin{columns}[t]
\begin{column}{0.55\textwidth}
\begin{abox}[TP1 : \file{notebooks/TP1_noyaux_convolution.ipynb}]
\begin{itemize}
  \item Charger une image réelle et la manipuler comme une matrice.
  \item Construire les noyaux classiques : flou, netteté, contours, relief.
  \item Appliquer un noyau et \textbf{interpréter} la carte d'activation.
  \item Refaire un calcul de convolution \textbf{à la main} sur une matrice
  $5 \times 5$, puis vérifier avec \file{scipy}.
\end{itemize}
\medskip
Durée conseillée : \textbf{1 h 30} --- barème : \textbf{10 points}.
\end{abox}
\end{column}
\begin{column}{0.42\textwidth}
\begin{wbox}[Avant de commencer]
Ne codez pas avant d'avoir répondu par écrit : « que doit produire ce noyau ? ».
L'intuition précède le code ; c'est la même démarche que dans un projet réel.
\end{wbox}
\begin{solbox}[Ce qui sera évalué]
La cohérence entre votre prédiction et le résultat observé, et la qualité de
l'interprétation des cartes obtenues.
\end{solbox}
\end{column}
\end{columns}
\end{frame}

\begin{frame}{Ce que le chapitre 2 va préciser}
\begin{columns}[t]
\begin{column}{0.48\textwidth}
\begin{kbox}[Trois questions ouvertes]
\begin{itemize}
  \item Comment se calcule \emph{exactement} une convolution, élément par élément ?
  \item Que deviennent les bords de l'image ? (\textbf{padding})
  \item Peut-on balayer plus vite ? (\textbf{stride})
\end{itemize}
\end{kbox}
\end{column}
\begin{column}{0.48\textwidth}
\begin{kbox}[Et une question de taille]
Comment prédire la \textbf{dimension} de la sortie d'une couche convolutive
sans la calculer ?

\medskip
C'est la formule du chapitre 2 --- et l'erreur la plus fréquente en
pratique.
\end{kbox}
\end{column}
\end{columns}
\end{frame}
